\documentclass[11pt]{article}
\usepackage{../common/ai204}
\renewcommand{\doctitle}{Assignments}
\begin{document}
\aicover{Assignment Booklet}{5 unit-wise assignments \textbullet\ think-out-of-the-box questions \textbullet\ open-ended coding problems with starter code}{Creativity \textbullet\ Critical thinking \textbullet\ Implementation \textbullet\ Responsible AI}
\section*{Instructions}
\begin{itemize}\item Each assignment has \textbf{Part A} (4 think-out-of-the-box questions, answer in 250--400 words each) and \textbf{Part B} (3 open-ended coding problems).
\item Starter code for Part B is in \texttt{09\_Assignments/notebooks/}; it runs as-is and prints a baseline. Complete every \texttt{TODO}.
\item There is no single correct answer. Marks reward originality, sound reasoning, working code, clear plots and honest discussion of limitations.
\item Submit a PDF report plus a GitHub link. You may use AI tools only if you state how you used them; copied or undisclosed AI text earns zero.\end{itemize}
\begin{center}\begin{tabular}{l p{9.5cm} c}\rowcolor{navy}\textcolor{gold}{\textbf{Criterion}}&\textcolor{gold}{\textbf{What earns full marks}}&\textcolor{gold}{\textbf{Marks}}\\
Originality&Ideas not found in the textbook; creative but defensible&5\\
\rowcolor{ivory}Reasoning&Correct use of course concepts; arguments with evidence&5\\
Implementation&Working, readable code; all TODOs completed; tests or checks&6\\
\rowcolor{ivory}Analysis&Plots, comparisons, honest limitations&4\\\midrule
\textbf{Total}&&\textbf{20}\\\end{tabular}\end{center}
\begin{center}\begin{tabular}{c l c c}\rowcolor{navy}\textcolor{gold}{\textbf{No.}}&\textcolor{gold}{\textbf{Title}}&\textcolor{gold}{\textbf{Unit / CO}}&\textcolor{gold}{\textbf{Due}}\\
1&Agents in the Wild&Unit I / CO204.1&End of Week 3\\
\rowcolor{ivory}2&Search Beyond the Textbook&Unit II / CO204.2&End of Week 6\\
3&Logic Meets the Real World&Unit III / CO204.3&End of Week 9\\
\rowcolor{ivory}4&Reasoning When Nothing Is Certain&Unit IV / CO204.4&End of Week 12\\
5&Learning Machines and Responsible AI&Unit V / CO204.5&End of Week 15\\
\end{tabular}\end{center}
\newpage
\unitbanner{I}{Assignment 1: Agents in the Wild}{CO204.1 \textbullet\ Due: End of Week 3}
\subsection*{\color{navy}Part A --- Think out of the box}
\begin{tcolorbox}[enhanced,colback=goldL,colframe=gold!80!black,arc=2mm,boxrule=0.8pt,left=3mm]
\textbf{\color{navy}A1.1}\quad If a self-driving car is perfectly rational for its performance measure but still hits a pedestrian in a sudden event, is the car 'to blame'? Argue from the definition of rationality, then propose how the performance measure itself should change.\hfill\btl{L5}
\end{tcolorbox}
\begin{tcolorbox}[enhanced,colback=goldL,colframe=gold!80!black,arc=2mm,boxrule=0.8pt,left=3mm]
\textbf{\color{navy}A1.2}\quad Invent a new agent type that does not fit any of the five classical architectures (reflex, model, goal, utility, learning). Describe its internal structure with a diagram and give a real problem where it beats all five.\hfill\btl{L6}
\end{tcolorbox}
\begin{tcolorbox}[enhanced,colback=goldL,colframe=gold!80!black,arc=2mm,boxrule=0.8pt,left=3mm]
\textbf{\color{navy}A1.3}\quad Your phone's keyboard autocomplete is an agent. Write its PEAS, classify its environment along all seven properties, and identify which property makes it hardest to design.\hfill\btl{L4}
\end{tcolorbox}
\begin{tcolorbox}[enhanced,colback=goldL,colframe=gold!80!black,arc=2mm,boxrule=0.8pt,left=3mm]
\textbf{\color{navy}A1.4}\quad Design a 'Reverse Turing Test' in which a machine must decide whether it is talking to a human or another AI. What questions would it ask, and why might this become important by 2030?\hfill\btl{L5}
\end{tcolorbox}
\subsection*{\color{navy}Part B --- Open-ended coding problems}
\subsubsection*{B1.1\quad Open-ended: a smarter vacuum world}
Extend the vacuum world to an n$\times$n grid where dirt re-appears randomly. Design the best agent you can and justify its performance measure. Compare at least three agent designs over 500 runs.
\begin{lstlisting}[style=py]
import random
def run(agent, n=4, steps=200, p_dirt=0.05, seed=0):
    rng = random.Random(seed); dirt = {(r, c): rng.random() < 0.5 for r in range(n) for c in range(n)}
    pos, score, mem = (0, 0), 0, {}
    for _ in range(steps):
        action = agent(pos, dirt[pos], mem, n)
        if action == "Suck" and dirt[pos]: dirt[pos] = False; score += 10
        elif action in ("Up", "Down", "Left", "Right"):
            dr, dc = {"Up": (-1, 0), "Down": (1, 0), "Left": (0, -1), "Right": (0, 1)}[action]
            pos = (min(max(pos[0] + dr, 0), n - 1), min(max(pos[1] + dc, 0), n - 1)); score -= 1
        for k in dirt:                                   # dirt re-appears
            if rng.random() < p_dirt: dirt[k] = True
    return score

def random_agent(pos, dirty, mem, n):
    return "Suck" if dirty else random.choice(["Up", "Down", "Left", "Right"])

print("Random agent mean score:", sum(run(random_agent, seed=s) for s in range(50)) / 50)
# TODO 1: write a model-based agent that remembers when it last visited each cell.
# TODO 2: write a utility-based agent that moves towards the cell most likely to be dirty.
# TODO 3: plot the score distributions of all agents and explain the winner.
\end{lstlisting}
\noindent{\small\bfseries\color{gold!70!black}Baseline output of the starter code}
\begin{lstlisting}[style=out]
Random agent mean score: 439.1
\end{lstlisting}
\subsubsection*{B1.2\quad Open-ended: formulate your own search problem}
Pick a real problem from your daily life (bus route to campus, cooking a meal, packing a bag). Formulate it with the five components and solve a small instance with BFS. Report the state-space size.
\begin{lstlisting}[style=py]
from collections import deque
# Example formulation: the classic Farmer-Wolf-Goat-Cabbage river crossing.
# State: tuple of sides (0 = left, 1 = right) for (farmer, wolf, goat, cabbage)
def safe(s):
    f, w, g, c = s
    return not ((w == g != f) or (g == c != f))
def successors(s):
    f = s[0]
    for i in range(4):                     # farmer crosses alone (i=0) or with item i
        if s[i] == f:
            t = list(s); t[0] = 1 - f
            if i: t[i] = 1 - f
            t = tuple(t)
            if safe(t): yield t
def bfs(start, goal):
    q, par = deque([start]), {start: None}
    while q:
        s = q.popleft()
        if s == goal:
            path = []
            while s: path.append(s); s = par[s]
            return path[::-1]
        for n in successors(s):
            if n not in par: par[n] = s; q.append(n)
path = bfs((0, 0, 0, 0), (1, 1, 1, 1))
print(f"Solved in {len(path) - 1} crossings:", path)
# TODO: replace this with YOUR problem. State clearly: states, initial state, actions,
#       transition model, goal test and path cost. Count reachable states.
\end{lstlisting}
\noindent{\small\bfseries\color{gold!70!black}Baseline output of the starter code}
\begin{lstlisting}[style=out]
Solved in 7 crossings: [(0, 0, 0, 0), (1, 0, 1, 0), (0, 0, 1, 0), (1, 1, 1, 0), (0, 1, 0, 0), (1, 1, 0, 1), (0, 1, 0, 1), (1, 1, 1, 1)]
\end{lstlisting}
\subsubsection*{B1.3\quad Open-ended: performance-measure stress test}
Write two performance measures for the same agent that lead to visibly different behaviour. Show the difference with a plot and explain which you would ship.
\begin{lstlisting}[style=py]
# Starter idea: an elevator agent choosing which floor to serve next.
import random
def simulate(choose, steps=300, floors=10, seed=1):
    rng = random.Random(seed); waiting, pos, waits = {}, 0, []
    for t in range(steps):
        if rng.random() < 0.6: waiting.setdefault(rng.randrange(floors), []).append(t)
        if waiting:
            target = choose(pos, waiting)
            pos += (target > pos) - (target < pos)          # move one floor per time step
            if pos in waiting: waits += [t - a for a in waiting.pop(pos)]
    return sum(waits) / max(len(waits), 1), max(waits or [0])
nearest = lambda pos, w: min(w, key=lambda f: abs(f - pos))
oldest  = lambda pos, w: min(w, key=lambda f: min(w[f]))
print("nearest-first (mean wait, max wait):", simulate(nearest))
print("oldest-first  (mean wait, max wait):", simulate(oldest))
# TODO: define the performance measure you think is fair, design an agent for it,
#       and discuss the trade-off between mean and worst-case waiting time.
\end{lstlisting}
\noindent{\small\bfseries\color{gold!70!black}Baseline output of the starter code}
\begin{lstlisting}[style=out]
nearest-first (mean wait, max wait): (4.331460674157303, 20)
oldest-first  (mean wait, max wait): (4.936781609195402, 16)
\end{lstlisting}
\newpage
\unitbanner{II}{Assignment 2: Search Beyond the Textbook}{CO204.2 \textbullet\ Due: End of Week 6}
\subsection*{\color{navy}Part A --- Think out of the box}
\begin{tcolorbox}[enhanced,colback=goldL,colframe=gold!80!black,arc=2mm,boxrule=0.8pt,left=3mm]
\textbf{\color{navy}A2.1}\quad A* is optimal with an admissible heuristic. Give a real situation where you would deliberately use an inadmissible heuristic, and argue why the loss of optimality is acceptable.\hfill\btl{L5}
\end{tcolorbox}
\begin{tcolorbox}[enhanced,colback=goldL,colframe=gold!80!black,arc=2mm,boxrule=0.8pt,left=3mm]
\textbf{\color{navy}A2.2}\quad Design a heuristic for finding the fastest route between two classrooms on your campus that accounts for stairs, lifts and crowds at bell time. Is it admissible? Prove or disprove it.\hfill\btl{L6}
\end{tcolorbox}
\begin{tcolorbox}[enhanced,colback=goldL,colframe=gold!80!black,arc=2mm,boxrule=0.8pt,left=3mm]
\textbf{\color{navy}A2.3}\quad Why can a chess engine look 20 moves ahead but a cricket-captain AI cannot? Analyse the difference in branching factor, observability and stochasticity.\hfill\btl{L4}
\end{tcolorbox}
\begin{tcolorbox}[enhanced,colback=goldL,colframe=gold!80!black,arc=2mm,boxrule=0.8pt,left=3mm]
\textbf{\color{navy}A2.4}\quad Invent a two-player game that is small enough to solve with Minimax but would be interesting for humans. Define its rules and predict who wins with perfect play.\hfill\btl{L6}
\end{tcolorbox}
\subsection*{\color{navy}Part B --- Open-ended coding problems}
\subsubsection*{B2.1\quad Open-ended: heuristic design challenge}
Invent at least two new admissible heuristics for the 8-puzzle (beyond misplaced tiles and Manhattan distance) and compare nodes expanded on 50 random solvable puzzles.
\begin{lstlisting}[style=py]
import heapq, itertools, random
GOAL = (1, 2, 3, 4, 5, 6, 7, 8, 0)
def succ(s):
    i = s.index(0); r, c = divmod(i, 3)
    for dr, dc in ((1, 0), (-1, 0), (0, 1), (0, -1)):
        nr, nc = r + dr, c + dc
        if 0 <= nr < 3 and 0 <= nc < 3:
            j = nr * 3 + nc; t = list(s); t[i], t[j] = t[j], t[i]; yield tuple(t)
def manhattan(s):
    return sum(abs(k // 3 - (v - 1) // 3) + abs(k % 3 - (v - 1) % 3) for k, v in enumerate(s) if v)
def astar(start, h):
    tie = itertools.count(); fr = [(h(start), 0, next(tie), start)]; g = {start: 0}; n = 0
    while fr:
        f, gs, _, s = heapq.heappop(fr); n += 1
        if s == GOAL: return gs, n
        for t in succ(s):
            if gs + 1 < g.get(t, 99): g[t] = gs + 1; heapq.heappush(fr, (gs + 1 + h(t), gs + 1, next(tie), t))
def scramble(k, seed):
    rnd = random.Random(seed); s = GOAL
    for _ in range(k): s = rnd.choice(list(succ(s)))
    return s
puzzles = [scramble(30, i) for i in range(10)]
print("Manhattan: mean nodes expanded =", sum(astar(p, manhattan)[1] for p in puzzles) / len(puzzles))
# TODO 1: implement 'linear conflict' (Manhattan + 2 per pair of tiles in the same row/col in reverse order).
# TODO 2: invent your own heuristic and prove it is admissible.
# TODO 3: run on 50 puzzles and plot nodes expanded for each heuristic.
\end{lstlisting}
\noindent{\small\bfseries\color{gold!70!black}Baseline output of the starter code}
\begin{lstlisting}[style=out]
Manhattan: mean nodes expanded = 49.7
\end{lstlisting}
\subsubsection*{B2.2\quad Open-ended: local search on a real optimisation problem}
Choose a real optimisation task (seating plan, delivery tour, playlist ordering) and solve it with hill climbing and simulated annealing. Report how often each reaches the best value you found.
\begin{lstlisting}[style=py]
import math, random
random.seed(0)
cities = [(random.random(), random.random()) for _ in range(25)]
def length(t): return sum(math.dist(cities[t[i]], cities[t[i - 1]]) for i in range(len(t)))
def two_opt_neighbour(t):
    i, j = sorted(random.sample(range(len(t)), 2)); return t[:i] + t[i:j][::-1] + t[j:]
def anneal(T=1.0, cool=0.995, steps=20000):
    cur = list(range(len(cities))); random.shuffle(cur); best = cur
    for _ in range(steps):
        nxt = two_opt_neighbour(cur); d = length(nxt) - length(cur)
        if d < 0 or random.random() < math.exp(-d / T): cur = nxt
        if length(cur) < length(best): best = cur
        T *= cool
    return length(best)
print("Simulated annealing tour length:", round(anneal(), 3))
# TODO: implement hill climbing with random restarts, compare over 20 seeds,
#       and try a different cooling schedule. Which matters more: schedule or neighbourhood?
\end{lstlisting}
\noindent{\small\bfseries\color{gold!70!black}Baseline output of the starter code}
\begin{lstlisting}[style=out]
Simulated annealing tour length: 4.184
\end{lstlisting}
\subsubsection*{B2.3\quad Open-ended: build and break a game AI}
Write a Minimax + alpha--beta agent for Connect-Three on a 4$\times$4 board (or your own game). Then write an evaluation function for a depth-limited version and show a position where it blunders.
\begin{lstlisting}[style=py]
import math
ROWS, COLS, WIN = 4, 4, 3
def moves(b): return [c for c in range(COLS) if b[0][c] == "."]
def drop(b, c, p):
    nb = [row[:] for row in b]
    for r in range(ROWS - 1, -1, -1):
        if nb[r][c] == ".": nb[r][c] = p; return nb
def winner(b):
    for r in range(ROWS):
        for c in range(COLS):
            for dr, dc in ((0, 1), (1, 0), (1, 1), (1, -1)):
                cells = [(r + k * dr, c + k * dc) for k in range(WIN)]
                if all(0 <= x < ROWS and 0 <= y < COLS for x, y in cells):
                    v = {b[x][y] for x, y in cells}
                    if len(v) == 1 and "." not in v: return v.pop()
    return None
def ab(b, p, a=-math.inf, be=math.inf):
    w = winner(b)
    if w: return 1 if w == "X" else -1
    if not moves(b): return 0
    best = -math.inf if p == "X" else math.inf
    for c in moves(b):
        v = ab(drop(b, c, p), "O" if p == "X" else "X", a, be)
        if p == "X": best = max(best, v); a = max(a, v)
        else: best = min(best, v); be = min(be, v)
        if a >= be: break
    return best
empty = [["."] * COLS for _ in range(ROWS)]
print("Value of the empty 4x4 Connect-Three board for X:", ab(empty, "X"))
# TODO: add a depth limit + evaluation function, measure nodes searched,
#       and find a position where the depth-limited agent makes a wrong move.
\end{lstlisting}
\noindent{\small\bfseries\color{gold!70!black}Baseline output of the starter code}
\begin{lstlisting}[style=out]
Value of the empty 4x4 Connect-Three board for X: 1
\end{lstlisting}
\newpage
\unitbanner{III}{Assignment 3: Logic Meets the Real World}{CO204.3 \textbullet\ Due: End of Week 9}
\subsection*{\color{navy}Part A --- Think out of the box}
\begin{tcolorbox}[enhanced,colback=goldL,colframe=gold!80!black,arc=2mm,boxrule=0.8pt,left=3mm]
\textbf{\color{navy}A3.1}\quad Translate the Sanjivani attendance rule (minimum attendance required to sit the exam, with medical exceptions) into first-order logic. Which parts resist precise logical encoding, and why?\hfill\btl{L4}
\end{tcolorbox}
\begin{tcolorbox}[enhanced,colback=goldL,colframe=gold!80!black,arc=2mm,boxrule=0.8pt,left=3mm]
\textbf{\color{navy}A3.2}\quad LLMs often answer logic puzzles wrongly. Is a resolution theorem prover 'smarter' than an LLM? Argue both sides with a concrete example.\hfill\btl{L5}
\end{tcolorbox}
\begin{tcolorbox}[enhanced,colback=goldL,colframe=gold!80!black,arc=2mm,boxrule=0.8pt,left=3mm]
\textbf{\color{navy}A3.3}\quad Design a knowledge graph of your own life (people, places, courses, skills). What new facts can you infer with three rules? What privacy risks would publishing it create?\hfill\btl{L6}
\end{tcolorbox}
\begin{tcolorbox}[enhanced,colback=goldL,colframe=gold!80!black,arc=2mm,boxrule=0.8pt,left=3mm]
\textbf{\color{navy}A3.4}\quad A semantic network says 'penguins are birds' and 'birds fly'. Propose a principled way to handle exceptions without rewriting every rule, and evaluate its weaknesses.\hfill\btl{L5}
\end{tcolorbox}
\subsection*{\color{navy}Part B --- Open-ended coding problems}
\subsubsection*{B3.1\quad Open-ended: your own expert knowledge base}
Encode at least 15 facts and 8 rules about a domain you know well (cricket rules, a recipe, a game). Answer five non-trivial queries with forward chaining and print the derivation trace.
\begin{lstlisting}[style=py]
FACTS = {("course", "AI"), ("course", "DSA"), ("prereq", "DSA", "AI"), ("passed", "asha", "DSA")}
RULES = [  # (head, [body]) with variables starting with '?'
    (("can_register", "?s", "?c"), [("passed", "?s", "?p"), ("prereq", "?p", "?c")]),
]
def match(pat, fact, th):
    if len(pat) != len(fact): return None
    th = dict(th)
    for p, f in zip(pat, fact):
        p = th.get(p, p)
        if p.startswith("?"): th[p] = f
        elif p != f: return None
    return th
def forward(facts):
    facts = set(facts); changed = True
    while changed:
        changed = False
        for head, body in RULES:
            thetas = [{}]
            for lit in body:
                thetas = [t2 for t in thetas for f in facts if (t2 := match(lit, f, t)) is not None]
            for t in thetas:
                new = tuple(t.get(x, x) for x in head)
                if new not in facts: facts.add(new); changed = True; print("derived", new)
    return facts
forward(FACTS)
# TODO: replace with YOUR domain (>= 15 facts, >= 8 rules, including one recursive rule),
#       add an explanation trace for each derived fact, and pose five queries.
\end{lstlisting}
\noindent{\small\bfseries\color{gold!70!black}Baseline output of the starter code}
\begin{lstlisting}[style=out]
derived ('can_register', 'asha', 'AI')
\end{lstlisting}
\subsubsection*{B3.2\quad Open-ended: resolution vs model checking}
Implement CNF conversion for propositional sentences and compare the speed of resolution and truth-table model checking as the number of symbols grows from 4 to 16.
\begin{lstlisting}[style=py]
import itertools, random, time
def random_kb(n, m, seed=0):
    rnd = random.Random(seed)
    return [[(rnd.randrange(n), rnd.random() < 0.5) for _ in range(3)] for _ in range(m)]
def satisfiable_tt(kb, n):
    for vals in itertools.product([False, True], repeat=n):
        if all(any(vals[v] == sign for v, sign in cl) for cl in kb): return True
    return False
for n in (4, 8, 12, 16):
    kb = random_kb(n, int(4.2 * n)); t = time.time()
    print(f"n={n:2d}: satisfiable={satisfiable_tt(kb, n)}  truth-table time={time.time() - t:.3f}s")
# TODO 1: implement propositional resolution (or DPLL) and time it on the same KBs.
# TODO 2: plot time vs n on a log scale and explain the curve. What is special about m/n ~ 4.2?
\end{lstlisting}
\noindent{\small\bfseries\color{gold!70!black}Baseline output of the starter code}
\begin{lstlisting}[style=out]
n= 4: satisfiable=True  truth-table time=0.000s
n= 8: satisfiable=True  truth-table time=0.000s
n=12: satisfiable=False  truth-table time=0.011s
n=16: satisfiable=True  truth-table time=0.043s
\end{lstlisting}
\subsubsection*{B3.3\quad Open-ended: mini ontology with SPARQL}
Build a small ontology ($\ge$ 20 triples) for the AI\&DS department and answer four questions with SPARQL queries.
\begin{lstlisting}[style=py]
try:
    from rdflib import Graph, Namespace, RDF, Literal
    EX = Namespace("http://example.org/aids#"); g = Graph()
    g.add((EX.AI, RDF.type, EX.Course)); g.add((EX.AI, EX.hasCredits, Literal(3)))
    g.add((EX.AI, EX.taughtIn, EX.Semester3)); g.add((EX.DSA, EX.prerequisiteOf, EX.AI))
    q = "SELECT ?c WHERE { ?c <http://example.org/aids#prerequisiteOf> <http://example.org/aids#AI> }"
    print("Prerequisites of AI:", [str(r.c).split("#")[1] for r in g.query(q)])
except ImportError:
    print("Install rdflib first:  pip install rdflib")
# TODO: model courses, faculty, labs and skills (>= 20 triples) and write four SPARQL queries,
#       one of which needs a property path (e.g. prerequisites of prerequisites).
\end{lstlisting}
\noindent{\small\bfseries\color{gold!70!black}Baseline output of the starter code}
\begin{lstlisting}[style=out]
Prerequisites of AI: ['DSA']
\end{lstlisting}
\newpage
\unitbanner{IV}{Assignment 4: Reasoning When Nothing Is Certain}{CO204.4 \textbullet\ Due: End of Week 12}
\subsection*{\color{navy}Part A --- Think out of the box}
\begin{tcolorbox}[enhanced,colback=goldL,colframe=gold!80!black,arc=2mm,boxrule=0.8pt,left=3mm]
\textbf{\color{navy}A4.1}\quad A COVID-style rapid test is advertised as '95\% accurate'. Explain to your family, without formulas, why a positive result might still be more likely wrong than right. When does it become reliable?\hfill\btl{L5}
\end{tcolorbox}
\begin{tcolorbox}[enhanced,colback=goldL,colframe=gold!80!black,arc=2mm,boxrule=0.8pt,left=3mm]
\textbf{\color{navy}A4.2}\quad Design a STRIPS model for 'getting ready for college in the morning'. Where does the classical-planning assumption fail in real life, and how would you repair the plan on the fly?\hfill\btl{L6}
\end{tcolorbox}
\begin{tcolorbox}[enhanced,colback=goldL,colframe=gold!80!black,arc=2mm,boxrule=0.8pt,left=3mm]
\textbf{\color{navy}A4.3}\quad Is 'tall' better modelled by probability or fuzzy logic? Construct a situation where using the wrong one leads to a bad decision.\hfill\btl{L4}
\end{tcolorbox}
\begin{tcolorbox}[enhanced,colback=goldL,colframe=gold!80!black,arc=2mm,boxrule=0.8pt,left=3mm]
\textbf{\color{navy}A4.4}\quad Bayesian networks need a causal structure. What goes wrong if an expert draws an arrow in the wrong direction? Propose a test to detect it from data.\hfill\btl{L5}
\end{tcolorbox}
\subsection*{\color{navy}Part B --- Open-ended coding problems}
\subsubsection*{B4.1\quad Open-ended: build a Bayesian network for a decision you care about}
Choose a diagnosis or decision (laptop won't boot, will it rain at college, should I take the bus). Build a network with $\ge$ 5 nodes, justify every CPT, and answer three queries with enumeration.
\begin{lstlisting}[style=py]
import itertools
# Starter: Rain -> WetRoad <- Sprinkler ; WetRoad -> Slip
P_R, P_S = 0.3, 0.2
P_W = {(True, True): 0.99, (True, False): 0.9, (False, True): 0.8, (False, False): 0.05}
P_SL = {True: 0.2, False: 0.01}
def joint(r, s, w, sl):
    p = (P_R if r else 1 - P_R) * (P_S if s else 1 - P_S)
    p *= P_W[(r, s)] if w else 1 - P_W[(r, s)]
    return p * (P_SL[w] if sl else 1 - P_SL[w])
def query(var_index, evidence):
    num = {True: 0.0, False: 0.0}
    for vals in itertools.product((True, False), repeat=4):
        if all(vals[i] == v for i, v in evidence.items()): num[vals[var_index]] += joint(*vals)
    z = sum(num.values()); return {k: round(v / z, 4) for k, v in num.items()}
print("P(Rain | Slip) =", query(0, {3: True}))
# TODO: design YOUR network (>= 5 nodes), justify CPTs from data or experts,
#       answer three queries and run a sensitivity analysis on one CPT entry.
\end{lstlisting}
\noindent{\small\bfseries\color{gold!70!black}Baseline output of the starter code}
\begin{lstlisting}[style=out]
P(Rain | Slip) = {True: 0.6222, False: 0.3778}
\end{lstlisting}
\subsubsection*{B4.2\quad Open-ended: fuzzy controller you can feel}
Design a fuzzy controller for something physical (fan speed, phone brightness, shower temperature). Plot membership functions and the control surface, then argue whether a simple formula would do better.
\begin{lstlisting}[style=py]
import numpy as np
def tri(x, a, b, c): return np.maximum(np.minimum((x - a) / (b - a + 1e-9), (c - x) / (c - b + 1e-9)), 0)
def fan_speed(temp):
    cold, warm, hot = tri(temp, 0, 10, 22), tri(temp, 15, 25, 35), tri(temp, 28, 40, 52)
    z = np.linspace(0, 100, 501)
    low, med, high = tri(z, -1, 0, 50), tri(z, 20, 50, 80), tri(z, 50, 100, 101)
    agg = np.maximum.reduce([np.minimum(cold, low), np.minimum(warm, med), np.minimum(hot, high)])
    return float((agg * z).sum() / (agg.sum() + 1e-9))      # centroid defuzzification
for t in (12, 24, 31, 38): print(f"{t}  degC -> fan speed {fan_speed(t):.1f}%")
# TODO: add a second input (humidity), write >= 6 rules, plot the 3-D control surface,
#       and compare with a hand-written linear formula.
\end{lstlisting}
\noindent{\small\bfseries\color{gold!70!black}Baseline output of the starter code}
\begin{lstlisting}[style=out]
12  degC -> fan speed 17.0%
24  degC -> fan speed 50.0%
31  degC -> fan speed 59.0%
38  degC -> fan speed 83.0%
\end{lstlisting}
\subsubsection*{B4.3\quad Open-ended: planning with surprises}
Implement a forward STRIPS planner for a delivery robot. Then make one action randomly fail and add execution monitoring with re-planning. Measure how often goals are still reached.
\begin{lstlisting}[style=py]
from collections import deque
ACTIONS = {  # name: (preconditions, add, delete)
    "move(A,B)": ({"at(A)"}, {"at(B)"}, {"at(A)"}), "move(B,A)": ({"at(B)"}, {"at(A)"}, {"at(B)"}),
    "move(B,C)": ({"at(B)"}, {"at(C)"}, {"at(B)"}), "move(C,B)": ({"at(C)"}, {"at(B)"}, {"at(C)"}),
    "pick(A)": ({"at(A)", "parcel(A)", "free"}, {"holding"}, {"parcel(A)", "free"}),
    "drop(C)": ({"at(C)", "holding"}, {"parcel(C)", "free"}, {"holding"}),
}
def plan(state, goal):
    q, seen = deque([(frozenset(state), [])]), set()
    while q:
        s, p = q.popleft()
        if goal <= s: return p
        if s in seen: continue
        seen.add(s)
        for a, (pre, add, dele) in ACTIONS.items():
            if pre <= s: q.append(((s - dele) | add, p + [a]))
print(plan({"at(B)", "parcel(A)", "free"}, {"parcel(C)"}))
# TODO: make 'move' fail with probability 0.2, monitor the state after each step,
#       re-plan on failure, and report the success rate and extra steps over 1000 runs.
\end{lstlisting}
\noindent{\small\bfseries\color{gold!70!black}Baseline output of the starter code}
\begin{lstlisting}[style=out]
['move(B,A)', 'pick(A)', 'move(A,B)', 'move(B,C)', 'drop(C)']
\end{lstlisting}
\newpage
\unitbanner{V}{Assignment 5: Learning Machines and Responsible AI}{CO204.5 \textbullet\ Due: End of Week 15}
\subsection*{\color{navy}Part A --- Think out of the box}
\begin{tcolorbox}[enhanced,colback=goldL,colframe=gold!80!black,arc=2mm,boxrule=0.8pt,left=3mm]
\textbf{\color{navy}A5.1}\quad A decision tree for loan approval uses 'pin code' as its top split. Is the model biased even if it never sees religion or caste? Explain proxy discrimination and propose a test.\hfill\btl{L5}
\end{tcolorbox}
\begin{tcolorbox}[enhanced,colback=goldL,colframe=gold!80!black,arc=2mm,boxrule=0.8pt,left=3mm]
\textbf{\color{navy}A5.2}\quad Design an expert system that could replace your college's help-desk for one week. Which questions must still go to a human, and how would the system know?\hfill\btl{L6}
\end{tcolorbox}
\begin{tcolorbox}[enhanced,colback=goldL,colframe=gold!80!black,arc=2mm,boxrule=0.8pt,left=3mm]
\textbf{\color{navy}A5.3}\quad Why did expert systems decline in the 1990s while machine learning rose? Would a 2030 'expert system' built on LLMs suffer the same problems?\hfill\btl{L4}
\end{tcolorbox}
\begin{tcolorbox}[enhanced,colback=goldL,colframe=gold!80!black,arc=2mm,boxrule=0.8pt,left=3mm]
\textbf{\color{navy}A5.4}\quad Should AI-generated assignments be banned, allowed or required? Take a position, then write the strongest argument against it.\hfill\btl{L5}
\end{tcolorbox}
\subsection*{\color{navy}Part B --- Open-ended coding problems}
\subsubsection*{B5.1\quad Open-ended: interpretable model vs black box}
Train a decision tree and one other model on a public dataset. Compare accuracy, extract the tree's rules, and decide whether the accuracy gain of the other model is worth the loss of interpretability.
\begin{lstlisting}[style=py]
from sklearn.datasets import load_breast_cancer
from sklearn.model_selection import cross_val_score
from sklearn.tree import DecisionTreeClassifier
from sklearn.ensemble import RandomForestClassifier
X, y = load_breast_cancer(return_X_y=True)
for name, m in (("Decision tree (depth 3)", DecisionTreeClassifier(max_depth=3, random_state=0)),
                ("Random forest", RandomForestClassifier(n_estimators=200, random_state=0))):
    print(f"{name:24s} 5-fold accuracy = {cross_val_score(m, X, y, cv=5).mean():.3f}")
# TODO: print the tree's rules, compute recall for the malignant class, and write a
#       one-page recommendation for a hospital: which model, at what threshold, and why.
\end{lstlisting}
\noindent{\small\bfseries\color{gold!70!black}Baseline output of the starter code}
\begin{lstlisting}[style=out]
Decision tree (depth 3)  5-fold accuracy = 0.917
Random forest            5-fold accuracy = 0.960
\end{lstlisting}
\subsubsection*{B5.2\quad Open-ended: perceptron to multilayer network}
Implement a perceptron and a 2-layer network from scratch with NumPy. Show the perceptron failing on XOR and the network succeeding; visualise both decision boundaries.
\begin{lstlisting}[style=py]
import numpy as np
X = np.array([[0, 0], [0, 1], [1, 0], [1, 1]]); y_xor = np.array([0, 1, 1, 0])
w, b = np.zeros(2), 0.0
for epoch in range(20):
    for xi, yi in zip(X, y_xor):
        pred = int(w @ xi + b >= 0); w += (yi - pred) * xi; b += (yi - pred)
print("Perceptron predictions on XOR:", [int(w @ xi + b >= 0) for xi in X], "target:", y_xor.tolist())
# TODO 1: implement a 2-2-1 network with sigmoid units and back-propagation; train on XOR.
# TODO 2: plot decision boundaries for both models.
# TODO 3: what is the smallest network that solves 3-input parity? Find it experimentally.
\end{lstlisting}
\noindent{\small\bfseries\color{gold!70!black}Baseline output of the starter code}
\begin{lstlisting}[style=out]
Perceptron predictions on XOR: [1, 1, 0, 0] target: [0, 1, 1, 0]
\end{lstlisting}
\subsubsection*{B5.3\quad Open-ended: fairness audit}
Using a dataset with a sensitive attribute (or a synthetic one you generate), train a classifier, measure accuracy and selection rate per group, and test one mitigation technique.
\begin{lstlisting}[style=py]
import numpy as np
from sklearn.linear_model import LogisticRegression
rng = np.random.default_rng(0); n = 4000
group = rng.integers(0, 2, n)                         # synthetic sensitive attribute
skill = rng.normal(0, 1, n); proxy = skill + 0.8 * group + rng.normal(0, 0.5, n)
label = (skill + rng.normal(0, 0.5, n) > 0).astype(int)
Xf = np.c_[proxy]; model = LogisticRegression().fit(Xf, label); pred = model.predict(Xf)
for g in (0, 1):
    m = group == g
    print(f"group {g}: accuracy={np.mean(pred[m] == label[m]):.3f}  selection rate={pred[m].mean():.3f}")
# TODO: explain why the groups differ even though 'group' is not a feature,
#       apply one mitigation (re-weighting, group thresholds, removing the proxy),
#       and report the accuracy-fairness trade-off in a small model card.
\end{lstlisting}
\noindent{\small\bfseries\color{gold!70!black}Baseline output of the starter code}
\begin{lstlisting}[style=out]
group 0: accuracy=0.777  selection rate=0.343
group 1: accuracy=0.769  selection rate=0.646
\end{lstlisting}
\newpage
\end{document}